% Chapter VI source-only decoder successor, 17 October 2026.
% LaTeX fragment; requires amsmath and amssymb, without a theorem package.
% Agent-originated derivation; no native execution or apparatus realization.

\section{A source-only decoder for the retained interacting clock}
\label{sec:chvi-source-decoder}

\paragraph{Provenance and receiving assumptions.}
The decoder and bounds below are agent-originated derivations by OpenAI
ChatGPT and Codex from the retained Chapter~I clock observable and its
Chapter~VI route interpretation. The source equations are in
\texttt{CHAPTER\_I\_MODEL.md}, Sections~2 and~5: the exact source-lane
restrictions, $Z=(D+iC)/d$, the reference expectation
$h(t)=Me^{-i\Omega t}$, and the actual bound $|z(t)-h(t)|\leq E(t)$.
The original signed coupling, both field polarizations, even source energy
and counterterms remain in the actual evolution that supplies $z$.
The ideal POVM at apparatus entry is source-only. During a finite physical
readout the same full Hamiltonian can continue to correlate source and
field; an implemented joint instrument must then approximate this ideal
POVM on the specified entry states. No direct field readout is prescribed.
Section~\ref{sec:chvisd-finite-readout} supplies that finite control and
pointer construction with its explicit drift allowance.

We retain the chapter's conditional route law: the complete marker starts
in its declared exact hard-eigenbranch/soft-vacuum preparation $\Psi_0$,
uncorrelated with the route label, and route $j$ supplies
$e^{-iH_0\eta_j}\Psi_0$. The source output is the partial trace over the
\emph{complete} field. No field fragment is measured by this decoder.

\subsection{The quadrature embedded in all sixteen source states}

Let $|q_a\rangle$ be the normalized ordered eigenvectors of the original
four-state source matrix $Q$. For one chosen lane $a<b$, abbreviate
$|ab\rangle=|q_a\rangle\otimes|q_b\rangle$ and define operators on
$\mathbb C^4\otimes\mathbb C^4$ by
\begin{equation}
\begin{aligned}
 P&=P^{\rm lane}_{ab}=|ab\rangle\langle ab|+|ba\rangle\langle ba|,\\
 Z_{ab}&=|ab\rangle\langle ab|-|ba\rangle\langle ba|,\\
 X_{ab}&=PSP,\qquad Y_{ab}=iSZ_{ab},
\end{aligned}
\label{eq:chvi-decoder-embedded-paulis}
\end{equation}
where $S$ is the original full source swap. On the ordered lane
$(|ab\rangle,|ba\rangle)$ these are $I,\sigma_z,\sigma_x,\sigma_y$;
$Z_{ab}$, $X_{ab}$ and $Y_{ab}$ vanish on the other fourteen source
states. In particular $Y_{ab}$ is Hermitian,
$Z_{ab}^2=Y_{ab}^2=P$, and $Z_{ab}Y_{ab}+Y_{ab}Z_{ab}=0$.
The complex clock observable embedded in this lane is
\begin{equation}
\begin{gathered}
Z^{\rm clk}_{ab}=Z_{ab}+iY_{ab}
   =\frac{P(D+iC)P}{d_{ab}},
 \\
 d_{ab}=q_a-q_b<0.
\end{gathered}
\label{eq:chvi-decoder-clock-operator}
\end{equation}
Every displayed source operator acts as its tensor product with the
identity on the field. Thus this is the original source observable,
not an observable requiring joint access to a dressed hard-field state.

The projector $P$ commutes with $Q_L,Q_R,S,D,C$ and preserves the common
energy form domain. It consequently reduces the retained form-sum
Hamiltonian. The exact lane preparation stays in this lane under $H_0$.
Writing its reduced source state at age $t$ as $\rho_S(t)$, we therefore
have
\begin{equation}
\begin{aligned}
 z(t)&=\operatorname{Tr}[\rho_S(t)Z^{\rm clk}_{ab}],\\
 O_\chi&=\frac{e^{-i\chi}Z^{\rm clk}_{ab}
                    +e^{i\chi}(Z^{\rm clk}_{ab})^*}{2}
       =\cos\chi\,Z_{ab}+\sin\chi\,Y_{ab},\\
 O_\chi^2&=P,\qquad \|O_\chi\|=1,\qquad
 m_j(\chi)=\operatorname{Re}[e^{-i\chi}z(\eta_j)].
\end{aligned}
\label{eq:chvi-decoder-quadrature}
\end{equation}

For an explicitly complete binary measurement, assign every state outside
the nominated lane to label zero:
\begin{equation}
\begin{gathered}
E_1(\chi)=\frac{P-O_\chi}{2},\\
 E_0(\chi)=I_{16}-E_1(\chi).
\end{gathered}
\label{eq:chvi-decoder-complete-povm}
\end{equation}
These are orthogonal projections and sum to $I_{16}$. On the lane they
are $(P\mp O_\chi)/2$; on its complement $E_0=I$ and $E_1=0$.
The full binary observable is accordingly
$E_0-E_1=(I_{16}-P)+O_\chi$, rather than $O_\chi$ on the complement.
No population has been dropped, and no success-conditioned
renormalization is used. For lane-supported states the conditional
correct-label probabilities are
\begin{equation}
\begin{gathered}
p(0|0)=\frac{1+m_0(\chi)}{2},\\
 p(1|1)=\frac{1-m_1(\chi)}{2}.
\end{gathered}
\label{eq:chvi-decoder-correct-probabilities}
\end{equation}
These probabilities and their worst-label bounds do not require balanced
label priors.

\subsection{General ages and a source-only discrimination witness}

\noindent\textbf{Proposition (quadrature contrast).}
For two exact lane-supported source outputs, let
$z_j=z(\eta_j)$ and
$D_S=\tfrac12\|\rho_S(\eta_0)-\rho_S(\eta_1)\|_1$.
Then
\begin{equation}
 D_S\geq\frac{|z_0-z_1|}{2}.
\label{eq:chvi-decoder-source-distance}
\end{equation}
If $z_0\ne z_1$, the quadrature maximizing $m_0-m_1$ is
$\chi=\arg(z_0-z_1)$ modulo $2\pi$. This maximizes contrast within the
stated quadrature family; it is not an assertion that this family is the
unrestricted Helstrom measurement or that it maximizes worst-label
success for arbitrary unequal-radius outputs.

\emph{Proof.} For any real $\chi$, trace/operator-norm duality gives
$|m_0-m_1|\leq2D_S\|O_\chi\|=2D_S$.
Maximizing $\operatorname{Re}[e^{-i\chi}(z_0-z_1)]$ gives
$|z_0-z_1|$. On lane-supported inputs the completed binary measurement
also attains the corresponding average success
$1/2+|z_0-z_1|/4$ for balanced priors. $\square$

For the reference expectations $h_j=Me^{-i\Omega\eta_j}$ with the
retained $M>0$, put
$\bar\eta=(\eta_0+\eta_1)/2$ and
$v=\sin[\Omega(\eta_1-\eta_0)/2]$.
When $v\ne0$, choose
\begin{equation}
 \chi_* =\arg(h_0-h_1)
 =\frac\pi2-\Omega\bar\eta+
 \begin{cases}0,&v>0,\\ \pi,&v<0,\end{cases}
 \quad \pmod{2\pi}.
\label{eq:chvi-decoder-reference-axis}
\end{equation}
The reference quadrature means are then $+M|v|$ and $-M|v|$.
For a predeclared calibrated implementation of this axis, the retained
actual-clock bound implies
\begin{equation}
\begin{gathered}
\min_jp(j|j)\geq
 \frac{1+M|v|-\max_jE(\eta_j)}2,
 \\
 D_S\geq\max\left\{0,M|v|-\frac{E(\eta_0)+E(\eta_1)}2\right\}.
\end{gathered}
\label{eq:chvi-decoder-general-bounds}
\end{equation}
The second bound follows directly from
$|z_0-z_1|\geq|h_0-h_1|-E(\eta_0)-E(\eta_1)$.
If the reference contrast vanishes, the displayed distance lower bound
is zero and no orientation is selected by this contrast criterion.
Choosing an axis from reference data before a trial does not presume
knowledge of its unknown label. Parameter uncertainty in that chosen
axis must be included in its calibration error.

\subsection{A fixed-axis, quarter-cycle protocol}

The following protocol avoids choosing its nominal axis from an absolute
hard energy or from the measured trial output. Use the retained chart
\begin{equation}
\begin{gathered}
T_0=\frac{2\pi}{\Delta},\\
 \eta_0=\frac{T_0}{4},\\
 \eta_1=\frac{3T_0}{4},\\
 \chi=-\frac\pi2,
 \\
 d=1-\frac{\Omega}{\Delta}.
\end{gathered}
\label{eq:chvi-decoder-fixed-protocol}
\end{equation}
The exact retained lane data satisfy $0<\Omega\leq\Delta$, so $d\geq0$.
The new symbol $d$ in Eq.~\eqref{eq:chvi-decoder-fixed-protocol} denotes
the relative gap deficit; it is distinct from the signed source
difference $d_{ab}$ in Eq.~\eqref{eq:chvi-decoder-clock-operator}.
The ages describe the two candidate marker histories at a common apparatus
entry, not two label-dependent readout triggers. The same rotation,
pointer test and outcome assignment are used for both hypotheses.
The measured lane quadrature is $O_{-\pi/2}=-Y_{ab}$.
Its reference means, including their signs, are
\begin{equation}
\begin{gathered}
m_{0,\rm ref}=M\cos(\pi d/2),\\
 m_{1,\rm ref}=-M\cos(3\pi d/2).
\end{gathered}
\label{eq:chvi-decoder-fixed-means}
\end{equation}
Thus the two \emph{correct-label margins} are positive cosines, whereas
the raw means have opposite signs. The retained comparison error gives
\begin{equation}\begin{gathered}
p(0|0)\geq\frac{1+M\cos(\pi d/2)-E(\eta_0)}2,
 \\
 p(1|1)\geq\frac{1+M\cos(3\pi d/2)-E(\eta_1)}2.
\label{eq:chvi-decoder-nominal-probabilities}
\end{gathered}\end{equation}

Let outward inputs satisfy
\begin{equation}\begin{gathered}
0\leq m_*\leq M\leq1,\\
 0\leq d\leq d_*,\\
 \pi\leq\pi_*,\\
 \Omega\leq\Omega_*,\\
 E(t)\leq E_*\\
(0\leq t\leq T_0).
\label{eq:chvi-decoder-outward-inputs}
\end{gathered}\end{equation}
Here $m_*$ is a lower bound for the retained $M$, not a rounded display.
For example $d_*\geq1-\Omega_{\rm lower}/\Delta_{\rm upper}$ is safe
with positive denominator and the stated $\Omega\leq\Delta$ premise.
$E_*$ can be supplied by the retained cycle enclosure of
$E(t)=\exp(4g^2v_s^{\rm up}t^2)-1$, which is nondecreasing on this
interval. The signed $g$ and the original split remain unchanged.

Set $\Theta_0=3\pi_*d_*/2$. Whenever $\Theta_0\leq\sqrt2$,
\begin{equation}
 \boxed{\begin{gathered}
 \epsilon_{\rm clock}:=\max_j[1-p(j|j)]
 \leq\frac{1-m_*[1-\Theta_0^2/2]+E_*}{2}.
 \end{gathered}}
\label{eq:chvi-decoder-nominal-error}
\end{equation}
Indeed $\cos x\geq1-x^2/2$ for real $x$. Both angular offsets in
Eq.~\eqref{eq:chvi-decoder-fixed-means} have modulus at most $\Theta_0$.
The condition makes the common lower cosine bound nonnegative, so
replacing $M$ by its lower bound preserves the inequality. Without
that condition the same substitution need not be valid; the universal
fallback is the trivial error bound one.

\subsection{Timing, axis and preparation robustness}

Allow actual ages $t_j=\eta_j+\xi_j$ with $|\xi_j|\leq a$ and
$0\leq t_j\leq T_0$, and let the actual quadrature axis be
$\chi=-\pi/2+\zeta$ with $|\zeta|\leq b$. A sufficient age-window
condition is $a\leq T_0/4$. These are dimensionless ages and angular
axis errors; no SI time calibration is introduced.
The signed reference means now obey
\begin{equation}
\begin{aligned}
 m_{0,\rm ref}&=M\cos(-\pi d/2+\Omega\xi_0+\zeta),\\
 -m_{1,\rm ref}&=M\cos(-3\pi d/2+\Omega\xi_1+\zeta).
\end{aligned}
\label{eq:chvi-decoder-jittered-means}
\end{equation}
Both correct-label angular errors are bounded by
\begin{equation}
 \Theta=\frac{3\pi_*d_*}{2}+\Omega_*a+b.
\label{eq:chvi-decoder-angle-budget}
\end{equation}
Consequently, for $\Theta\leq\sqrt2$,
\begin{equation}
 \max_j[1-p(j|j)]
 \leq\frac{1-m_*[1-\Theta^2/2]+E_*}{2}.
\label{eq:chvi-decoder-jitter-error}
\end{equation}
The timing and axis losses enter through a squared displacement from the
reference antipodal extrema. The square includes the gap-deficit and
cross terms; it is not a claim that each uncertainty contributes an
independent additive quadratic error. This bound is on a source
probability, so it needs no absolute hard mean energy $\bar E$ or
absolute route amplitude phase. Indeed a common scalar shift of the
generator multiplies each branch propagator by a scalar phase, which
cancels in $U_j\rho U_j^*$ and hence in its source measurement
probabilities. That invariance does not remove the scalar contribution
from an interferometric comparison of the route amplitudes themselves.

Suppose an imperfect complete marker preparation $\rho$ satisfies
\begin{equation}
 \tfrac12\|\rho-|\Psi_0\rangle\langle\Psi_0|\|_1
 \leq\delta_{\rm prep},
\label{eq:chvi-decoder-preparation-budget}
\end{equation}
with the same generator and declared ages. Unitary evolution and partial
trace do not increase this distance. Each completed-POVM outcome
probability therefore changes by at most $\delta_{\rm prep}$: for an
effect $0\leq F\leq I$, the trace-distance variational bound is
$|\operatorname{Tr}[F(\rho-\sigma)]|\leq T(\rho,\sigma)$.
It follows that
\begin{equation}
\begin{gathered}
\boxed{\begin{gathered}
 \epsilon_{\rm dec}\leq
 \min\left\{1,
 \frac{1-m_*[1-\Theta^2/2]+E_*}{2}+\delta_{\rm prep}
 \right\},\\
 \Theta\leq\sqrt2.
 \end{gathered}}
\end{gathered}
\label{eq:chvi-decoder-total-error}
\end{equation}
This includes leakage into the other fourteen source states and
source--field preparation errors covered by
Eq.~\eqref{eq:chvi-decoder-preparation-budget}. The penalty is
$\delta_{\rm prep}$ for an outcome probability, rather than the
$2\delta_{\rm prep}$ bound for a general unitary amplitude.
If an apparatus construction subsequently bounds an additional per-label
outcome-probability error by $\delta_{\rm read}$, it may be added to
the right side before clipping at one. That implementation error is not
established by the ideal measurement algebra alone. In particular, the
finite-control Kraus operators may act jointly on source and field even
when all external controls and the pointer couple only to the source.
Their approximation to the ideal entry POVM, rather than an exact
source--field factorization during the pulse, is the required estimate.

If Eq.~\eqref{eq:chvi-decoder-preparation-budget} is established only on a
heralded ready branch, its resulting error bound is conditional on that
branch. With per-label failure probability at most $f$, counting a
failure as unsuccessful gives the all-attempt bound
$f+(1-f)\epsilon$, where $\epsilon\leq1$ is the complete ready-branch
error including preparation and readout. The failure outcome remains in
the instrument; this formula does not assert an additional implemented
flag or silently discard unsuccessful attempts.

\subsection{Rotation sign and the complete binary readout}

For a general nominal axis let $R_\chi=e^{-i\chi S/2}$. The original
swap satisfies $S^2=I_{16}$ and commutes with $P$, while it anticommutes
with $Z_{ab}$ on the lane. Hence
\begin{equation}
\begin{gathered}
R_\chi^*Z_{ab}R_\chi
   =\cos\chi\,Z_{ab}+\sin\chi\,Y_{ab}=O_\chi,
 \\
 E_1(\chi)=R_\chi^*|ba\rangle\langle ba|R_\chi.
\end{gathered}
\label{eq:chvi-decoder-rotation}
\end{equation}
In particular
\begin{equation}
\begin{gathered}
R=e^{+i\pi S/4},\\
 R^*Z_{ab}R=-Y_{ab},\\
 E_1=R^*|ba\rangle\langle ba|R,\\
 E_0=I_{16}-E_1.
\end{gathered}
\label{eq:chvi-decoder-fixed-readout}
\end{equation}
Thus rotating by $R$ and testing for $|ba\rangle$ realizes the specified
\emph{ideal POVM}; every other outcome is assigned label zero.
Equivalently, a three-valued $Z_{ab}$ read after the rotation can be
coarse-grained by assigning its zero eigenspace to label zero.
An exchange-angle error $R=e^{i(\pi/4+\alpha)S}$ gives quadrature-axis
error $\zeta=-2\alpha$; a bound on $|\alpha|$ must therefore be doubled
when used as the axis budget $b$. Other errors caused by a concurrently
running drift or field interaction need their own control estimate.

For a separately declared leakage-calibration experiment one may instead
use the three effects
\begin{equation}
\begin{gathered}
R^*|ab\rangle\langle ab|R,\\
 R^*|ba\rangle\langle ba|R,\\
 I_{16}-P.
\end{gathered}
\end{equation}
The third effect is a failure/leakage outcome and the three sum to
$I_{16}$. This is a different stated output alphabet; its existence as
a POVM does not supply a second physical flag apparatus, and conditioning
away its third outcome would require a separate probability account.

The present section supplies the complete source-only POVM and its
worst-label error budget. Section~\ref{sec:chvisd-finite-readout} implements
its rotation and eigenbasis test with finite bounded controls and a
two-state pointer. The durable memory, causal receiving channel and
microscopic control actuator form the next apparatus join.

